\chapter{Conclusiones y Recomendaciones}
\label{sec:conclusiones}

\section{Conclusiones Técnicas}
A partir del relevamiento sonoro continuo efectuado con la aplicación Sound Analyzer v2.7 sobre los datos reales exportados
y el análisis conforme a la normativa vigente, se concluye lo siguiente:

\begin{enumerate}
  \item El Nivel Sonoro Continuo Equivalente final obtenido durante el ensayo de $306,69\s$ con el minitorno Dremel en
        sus dos regímenes operativos ($15\text{ kRPM}$ y $25\text{ kRPM}$) alcanza los $L_{Aeq} = 84,66\dBA$.
  \item El nivel registrado cumple formalmente con el valor límite permisible de $85,00\dBA$ prescripto por la Res. 295/03
        MTEySS para jornadas de $8\h$.
  \item La dosis diaria de ruido proyectada a las 8 horas de turno equivale al $92,5\%$, no superando el $100\%$ legal.
  \item El puesto de trabajo califica como \textbf{Conforme} en el protocolo oficial de la Resolución SRT 85/2012.
\end{enumerate}

\section{Recomendaciones Preventivas}
Si bien el puesto se encuentra dentro de los límites legales, al situarse en el $92,5\%$ del umbral máximo admisible se
recomiendan las siguientes acciones preventivas:

\begin{itemize}
  \item \textbf{Mantenimiento y Control de Vibraciones:} Realizar inspección periódica del rodamiento del minitorno y
        descartar mandriles o discos desbalanceados que incrementen el ruido por vibración mecánica.
  \item \textbf{Protección Personal (EPP):} Recomendar el uso de protectores auditivos de inserción en jornadas prolongadas
        de pulido intensivo a $25\text{ kRPM}$.
  \item \textbf{Vigilancia de la Salud:} Realizar audiometrías tonales periódicas anuales al personal técnico de reparación.
\end{itemize}
